\section{Extended Ablation Studies}
\label{sec:ablation}

We perform extended ablation studies on latent-space image generation. Specifically, we vary the error tolerance $\delta$, the number of resolution stages $S$, and the spectral transformation $T_\Phi$, and report wall-clock latency, speedup, and ImageReward in Table~\ref{tab:ablation-quan}, with qualitative comparisons in Figures~\ref{fig:ablation-T} to \ref{fig:ablation-s}.

\begin{table*}[t!]
    \centering
    \caption{\textbf{Ablation studies on the spectral transformation $T_\Phi$, error tolerance $\delta$, and number of resolution stages $S$.} Speedup is calculated relative to the 7-step FLUX.1-dev baseline wall-clock runtime on a single H100 GPU.}
    \label{tab:ablation-quan}
    \centering
    \begin{tiny}
    \begin{tabular}{cccccc}
        \toprule
        $T_\Phi$ & $\delta$ & $S$ & Latency (s) $\downarrow$ & Speedup (s) $\uparrow$ & ImageReward $\uparrow$ \\
        \midrule
        FLUX (7 steps) & - & - & 2.08 & 1.00$\times$ & 0.920 \\
        \midrule
        DCT & 0.001 & 2 & \textbf{2.08} & \textbf{1.00$\times$} & 1.032 \\
        DWT & 0.001 & 2 & \underline{2.08} & \textbf{1.00$\times$} & \textbf{1.038} \\
        FFT & 0.001 & 2 & 2.08 & \textbf{1.00$\times$} & \underline{1.036} \\
        \midrule
        DCT & 0.001 & 3 & \textbf{2.05} & \textbf{1.01$\times$} & \underline{1.034} \\
        DWT & 0.001 & 3 & \underline{2.06} & \textbf{1.01$\times$} & \underline{1.034} \\
        FFT & 0.001 & 3 & 2.06 & \textbf{1.01$\times$} & \textbf{1.040} \\
        \midrule
        DCT & 0.001 & 4 & \textbf{2.05} & \textbf{1.02$\times$} & \textbf{1.008} \\
        DWT & 0.001 & 4 & \underline{2.05} & \underline{1.01$\times$} & \underline{1.002} \\
        FFT & 0.001 & 4 & 2.06 & \underline{1.01$\times$} & 0.992 \\
        \midrule
        DCT & 0.002 & 2 & \textbf{2.07} & \textbf{1.01$\times$} & \underline{1.032} \\
        DWT & 0.002 & 2 & \underline{2.08} & \underline{1.00$\times$} & 1.031 \\
        FFT & 0.002 & 2 & 2.08 & \underline{1.00$\times$} & \textbf{1.036} \\
        \midrule
        DCT & 0.002 & 3 & \textbf{2.05} & \textbf{1.02$\times$} & \underline{1.034} \\
        DWT & 0.002 & 3 & \underline{2.06} & \underline{1.01$\times$} & \underline{1.034} \\
        FFT & 0.002 & 3 & 2.06 & \underline{1.01$\times$} & \textbf{1.040} \\
        \midrule
        DCT & 0.002 & 4 & \textbf{2.05} & \textbf{1.02$\times$} & \textbf{1.008} \\
        DWT & 0.002 & 4 & \underline{2.05} & \underline{1.01$\times$} & \underline{1.002} \\
        FFT & 0.002 & 4 & 2.06 & \underline{1.01$\times$} & 0.992 \\
        \midrule
        DCT & 0.005 & 2 & \textbf{1.91} & \textbf{1.09$\times$} & \underline{1.019} \\
        DWT & 0.005 & 2 & \underline{1.92} & \textbf{1.09$\times$} & \textbf{1.026} \\
        FFT & 0.005 & 2 & 1.92 & \underline{1.08$\times$} & 1.014 \\
        \midrule
        DCT & 0.005 & 3 & \textbf{1.87} & \textbf{1.12$\times$} & \underline{1.009} \\
        DWT & 0.005 & 3 & \underline{1.88} & \underline{1.11$\times$} & \underline{1.009} \\
        FFT & 0.005 & 3 & 1.88 & \underline{1.11$\times$} & \textbf{1.028} \\
        \midrule
        DCT & 0.005 & 4 & \textbf{1.87} & \textbf{1.12$\times$} & \underline{1.005} \\
        DWT & 0.005 & 4 & \underline{1.87} & \underline{1.11$\times$} & \textbf{1.008} \\
        FFT & 0.005 & 4 & 1.87 & \underline{1.11$\times$} & 0.994 \\
        \midrule
        DCT & 0.01 & 2 & \textbf{1.75} & \textbf{1.19$\times$} & \textbf{1.039} \\
        DWT & 0.01 & 2 & \underline{1.76} & \underline{1.18$\times$} & \underline{1.025} \\
        FFT & 0.01 & 2 & 1.76 & \underline{1.18$\times$} & 1.016 \\
        \midrule
        DCT & 0.01 & 3 & \textbf{1.71} & \textbf{1.22$\times$} & 1.015 \\
        DWT & 0.01 & 3 & \underline{1.72} & \underline{1.21$\times$} & \underline{1.018} \\
        FFT & 0.01 & 3 & 1.72 & \underline{1.21$\times$} & \textbf{1.030} \\
        \midrule
        DCT & 0.01 & 4 & \textbf{1.71} & \textbf{1.22$\times$} & \textbf{1.014} \\
        DWT & 0.01 & 4 & \underline{1.71} & \textbf{1.22$\times$} & \underline{1.006} \\
        FFT & 0.01 & 4 & 1.72 & \underline{1.21$\times$} & 0.996 \\
        \midrule
        DCT & 0.02 & 2 & \textbf{1.60} & \textbf{1.30$\times$} & \underline{1.009} \\
        DWT & 0.02 & 2 & \underline{1.61} & \textbf{1.30$\times$} & \textbf{1.010} \\
        FFT & 0.02 & 2 & \underline{1.61} & \textbf{1.30$\times$} & 1.002 \\
        \midrule
        DCT & 0.02 & 3 & \textbf{1.53} & \textbf{1.36$\times$} & \underline{0.967} \\
        DWT & 0.02 & 3 & \underline{1.53} & \textbf{1.36$\times$} & 0.966 \\
        FFT & 0.02 & 3 & 1.53 & \textbf{1.36$\times$} & \textbf{0.980} \\
        \midrule
        DCT & 0.02 & 4 & \textbf{1.53} & \textbf{1.36$\times$} & \textbf{0.964} \\
        DWT & 0.02 & 4 & \underline{1.53} & \textbf{1.36$\times$} & \underline{0.955} \\
        FFT & 0.02 & 4 & 1.53 & \textbf{1.36$\times$} & 0.949 \\
        \midrule
        DCT & 0.05 & 2 & \textbf{1.44} & \textbf{1.44$\times$} & \underline{0.989} \\
        DWT & 0.05 & 2 & \underline{1.45} & \textbf{1.44$\times$} & \textbf{0.991} \\
        FFT & 0.05 & 2 & 1.45 & \underline{1.43$\times$} & 0.972 \\
        \midrule
        DCT & 0.05 & 3 & \textbf{1.34} & \textbf{1.55$\times$} & \underline{0.936} \\
        DWT & 0.05 & 3 & \underline{1.34} & \textbf{1.55$\times$} & \textbf{0.937} \\
        FFT & 0.05 & 3 & 1.35 & \underline{1.54$\times$} & 0.933 \\
        \midrule
        DCT & 0.05 & 4 & \textbf{1.33} & \textbf{1.56$\times$} & \textbf{0.918} \\
        DWT & 0.05 & 4 & \underline{1.34} & \textbf{1.56$\times$} & \underline{0.902} \\
        FFT & 0.05 & 4 & 1.34 & \underline{1.55$\times$} & 0.895 \\
        \midrule
        DCT & 0.1 & 2 & \textbf{1.29} & \textbf{1.62$\times$} & \underline{0.966} \\
        DWT & 0.1 & 2 & \underline{1.29} & \textbf{1.62$\times$} & \textbf{0.977} \\
        FFT & 0.1 & 2 & 1.30 & \underline{1.61$\times$} & 0.944 \\
        \midrule
        DCT & 0.1 & 3 & \textbf{1.15} & \textbf{1.81$\times$} & \underline{0.880} \\
        DWT & 0.1 & 3 & 1.17 & 1.78$\times$ & \textbf{0.892} \\
        FFT & 0.1 & 3 & \underline{1.17} & \underline{1.79$\times$} & 0.877 \\
        \midrule
        DCT & 0.1 & 4 & \textbf{1.15} & \textbf{1.82$\times$} & \underline{0.864} \\
        DWT & 0.1 & 4 & \underline{1.15} & \underline{1.81$\times$} & \textbf{0.867} \\
        FFT & 0.1 & 4 & 1.16 & 1.80$\times$ & 0.841 \\
        \bottomrule
    \end{tabular}
    \end{tiny}
\end{table*}


\paragraph{Computational efficiency and scaling.} The inference cost of Spectral Progressive Diffusion is determined by the number of denoising steps and the token count processed at each resolution stage. For a fixed pretrained model and resolution schedule, the per-sample cost is independent of the dataset size, so generating or evaluating more samples scales linearly with the number of samples. The spectral transformation and noise-expansion operations add only lightweight per-transition overhead compared to the DiT forward passes. 

\paragraph{Spectral Transformation $T_\Phi$.} In Fig. \ref{fig:ablation-T}, we observe that DCT, DWT, and FFT achieve similar high-fidelity performance. We default to DCT for its real-valued computation and native support of non-power-of-two resolution scaling.

\paragraph{Error tolerance $\delta$.} In Fig. \ref{fig:ablation-delta}, we observe that increasing $\delta$ improves efficiency, but results in ghosting and halo artifacts near detailed edges.

\paragraph{Resolution stages $S$.} In Fig. \ref{fig:ablation-s}, we find that increasing $S$ leads to marginal speedup improvements and little image quality degradation. Since the primary speedup derives from bypassing native high-resolution steps, $S=2$ achieves the largest speedup while higher $S$ provides diminishing efficiency gains by further optimizing the already-inexpensive low-resolution trajectory.




\newpage
\begin{figure}
    \centering
    \includegraphics[width=\linewidth]{supp_figures/ablation_T.pdf}
    \caption{\textbf{Qualitative ablation on $T_\Phi$.} We see that all spectral transformations attain similar image quality. }
    \label{fig:ablation-T}
\end{figure}

\newpage
\begin{figure}
    \centering
    \includegraphics[width=\linewidth]{supp_figures/ablation_delta.pdf}
    \caption{\textbf{Qualitative ablation on $\delta$.} We observe that increasing $\delta$ improves efficiency, but results in ghosting and halo artifacts near detailed edges. }
    \label{fig:ablation-delta}
\end{figure}

\newpage
\begin{figure}
    \centering
    \includegraphics[width=\linewidth]{supp_figures/ablation_s.pdf}
    \caption{\textbf{Qualitative ablation on $S$.} We find that increasing $S$ leads to marginal speedup improvements and little image quality degradation.}
    \label{fig:ablation-s}
\end{figure}
